\SourceLeaf{247}
\Article{235}
Wenn die Form
\[
AXX+2BXY+CYY\ldots F
\]
in das Produkt der beiden Formen
\[
axx+2bxy+cyy\ldots f,\quad\text{und}\quad a'x'x'+2b'x'y'+c'y'y'\ldots f'
\]
durch eine Substitution der Art
\[
\begin{aligned}
X&=pxx'+p'xy'+p''yx'+p'''yy'\\
Y&=qxx'+q'xy'+q''yx'+q'''yy'
\end{aligned}
\]
übergeht (was wir im Folgenden der Kürze wegen stets so ausdrücken: Wenn $F$ in $ff'$ durch die Substitution $p,p',p'',p''';\ q,q',q'',q'''$ übergeht\HistoricalNote{*)}{Bei dieser Bezeichnung muss man also sowohl die Reihenfolge der Koeffizienten $p,p'$ usw. als auch die der Formen $f,f'$ genau beachten. Man sieht leicht: Wird die Reihenfolge der Formen $f,f'$ umgekehrt, sodass die erste zur zweiten wird, so müssen die Koeffizienten $p',q'$ mit $p'',q''$ vertauscht werden, während jeder der übrigen an seinem Platz bleibt.}), nennen wir die Form $F$ einfach in $ff'$ \textit{transformierbar}; ist diese Transformation außerdem so beschaffen, dass die sechs Zahlen
\[
\begin{gathered}pq'-qp',\quad pq''-qp'',\quad pq'''-qp''',\\p'q''-q'p'',\quad p'q'''-q'p''',\quad p''q'''-q''p'''\end{gathered}
\]
keinen gemeinsamen Teiler haben, so nennen wir die Form $F$ aus den Formen $f,f'$ \textit{zusammengesetzt}.

Wir beginnen diese Untersuchung mit der allgemeinsten Annahme, dass die Form $F$ in $ff'$ durch die Substitution $p,p',p'',p''';\ q,q',q'',q'''$ übergeht, und entwickeln ihre Folgen. Dieser Annahme sind offenbar die folgenden neun Gleichungen genau gleichwertig (\textit{d.h.}, sobald diese Gleichungen gelten, geht $F$ durch die genannte Substitution in $ff'$ über, und umgekehrt):
\begin{align*}
App+2Bpq+Cqq&=aa'\quad\ldotp\,\ldotp\,\ldotp\tag*{[1]}\\
Ap'p'+2Bp'q'+Cq'q'&=ac'\quad\ldotp\,\ldotp\,\ldotp\tag*{[2]}\\
Ap''p''+2Bp''q''+Cq''q''&=ca'\quad\ldotp\,\ldotp\,\ldotp\tag*{[3]}\\
Ap'''p'''+2Bp'''q'''+Cq'''q'''&=cc'\quad\ldotp\,\ldotp\,\ldotp\tag*{[4]}\\
App'+B(pq'+qp')+Cqq'&=ab'\quad\ldotp\,\ldotp\,\ldotp\tag*{[5]}\\
App''+B(pq''+qp'')+Cqq''&=ba'\quad\ldotp\,\ldotp\,\ldotp\tag*{[6]}\\
Ap'p'''+B(p'q'''+q'p''')+Cq'q'''&=bc'\quad\ldotp\,\ldotp\,\ldotp\tag*{[7]}\\
Ap''p'''+B(p''q'''+q''p''')+Cq''q'''&=cb'\quad\ldotp\,\ldotp\,\ldotp\tag*{[8]}\\
\begin{aligned}A(pp'''+p'p'')&+B(pq'''+qp'''+p'q''+q'p'')\\&+C(qq'''+q'q'')\end{aligned}&=2bb'\quad\ldotp\,\ldotp\,\ldotp\tag*{[9]}
\end{align*}

\SourceLeaf{248}
Die Determinanten der Formen $F,f,f'$ seien jeweils $D,d,d'$; die größten gemeinsamen Teiler der Zahlen $A,2B,C;\ a,2b,c;\ a',2b',c'$ seien jeweils $M,m,m'$ (alle positiv genommen). Weiter bestimmen wir sechs ganze Zahlen $\mathfrak{A},\mathfrak{B},\mathfrak{C},\mathfrak{A}',\mathfrak{B}',\mathfrak{C}'$ so, dass
\[
\mathfrak{A}a+2\mathfrak{B}b+\mathfrak{C}c=m,\qquad
\mathfrak{A}'a'+2\mathfrak{B}'b'+\mathfrak{C}'c'=m'
\]
gilt. Schließlich bezeichnen wir die Zahlen
\[
\begin{gathered}pq'-qp',\quad pq''-qp'',\quad pq'''-qp''',\\p'q''-q'p'',\quad p'q'''-q'p''',\quad p''q'''-q''p'''\end{gathered}
\]
der Reihe nach mit $P,Q,R,S,T,U$, und ihr positiv genommener größter gemeinsamer Teiler sei $=k$. --- Setzen wir nun
\begin{align*}
App'''+B(pq'''+qp''')+Cqq'''&=bb'+\Delta\quad\ldotp\,\ldotp\,\ldotp\,\ldotp\tag*{[10]}
\end{align*}
so folgt aus Gleichung 9
\begin{align*}
Ap'p''+B(p'q''+q'p'')+Cq'q''&=bb'-\Delta\quad\ldotp\,\ldotp\,\ldotp\,\ldotp\tag*{[11]}
\end{align*}

Aus diesen elf Gleichungen $1\ldots11$ leiten wir die folgenden neuen ab\HistoricalNote{*)}{Diese Gleichungen entstehen so: 12 aus $5.5-1.2$; 13 aus $5.9-1.7-2.6$; 14 aus $10.11-6.7$; 15 aus $5.8+5.8+10.10+11.11-1.4-2.3-6.7-6.7$; 16 aus $8.9-3.7-4.6$; 17 aus $8.8-3.4$. Die übrigen sechs leitet man ebenso her, wenn man die Gleichungen $2,5,7$ jeweils mit den Gleichungen $3,6,8$ vertauscht und die übrigen $1,4,9,10,11$ weiterhin an derselben Stelle lässt; so entsteht 18 aus $6.6-1.3$ usw.}:
\begin{align*}
DPP&=d'aa\quad\ldotp\,\ldotp\,\ldotp\,\ldotp\,\ldotp\,\ldotp\,\ldotp\,\ldotp\,\ldotp\,\ldotp\,\ldotp\,\ldotp\,\ldotp\tag*{[12]}\\
DP(R-S)&=2d'ab\quad\ldotp\,\ldotp\,\ldotp\,\ldotp\,\ldotp\,\ldotp\,\ldotp\,\ldotp\,\ldotp\,\ldotp\,\ldotp\,\ldotp\tag*{[13]}\\
DPU&=d'ac-(\Delta\Delta-dd')\quad\ldotp\,\ldotp\,\ldotp\,\ldotp\,\ldotp\,\ldotp\,\ldotp\tag*{[14]}\\
D(R-S)^2&=4d'bb+2(\Delta\Delta-dd')\quad\ldotp\,\ldotp\,\ldotp\,\ldotp\,\ldotp\tag*{[15]}\\
D(R-S)U&=2d'bc\quad\ldotp\,\ldotp\,\ldotp\,\ldotp\,\ldotp\,\ldotp\,\ldotp\,\ldotp\,\ldotp\,\ldotp\,\ldotp\,\ldotp\tag*{[16]}\\
DUU&=d'cc\quad\ldotp\,\ldotp\,\ldotp\,\ldotp\,\ldotp\,\ldotp\,\ldotp\,\ldotp\,\ldotp\,\ldotp\,\ldotp\,\ldotp\,\ldotp\tag*{[17]}\\
DQQ&=da'a'\quad\ldotp\,\ldotp\,\ldotp\,\ldotp\,\ldotp\,\ldotp\,\ldotp\,\ldotp\,\ldotp\,\ldotp\,\ldotp\,\ldotp\tag*{[18]}\\
DQ(R+S)&=2da'b'\quad\ldotp\,\ldotp\,\ldotp\,\ldotp\,\ldotp\,\ldotp\,\ldotp\,\ldotp\,\ldotp\,\ldotp\,\ldotp\tag*{[19]}\\
DQT&=da'c'-(\Delta\Delta-dd')\quad\ldotp\,\ldotp\,\ldotp\,\ldotp\,\ldotp\,\ldotp\tag*{[20]}\\
D(R+S)^2&=4db'b'+2(\Delta\Delta-dd')\quad\ldotp\,\ldotp\,\ldotp\,\ldotp\,\ldotp\tag*{[21]}\\
D(R+S)T&=2db'c'\quad\ldotp\,\ldotp\,\ldotp\,\ldotp\,\ldotp\,\ldotp\,\ldotp\,\ldotp\,\ldotp\,\ldotp\,\ldotp\,\ldotp\tag*{[22]}\\
DTT&=dc'c'\quad\ldotp\,\ldotp\,\ldotp\,\ldotp\,\ldotp\,\ldotp\,\ldotp\,\ldotp\,\ldotp\,\ldotp\,\ldotp\,\ldotp\,\ldotp\tag*{[23]}
\end{align*}

Daraus folgen wiederum diese beiden:
\SourceLeaf{249}
\[
\begin{aligned}
0&=2d'aa(\Delta\Delta-dd')\\
0&=(\Delta\Delta-dd')^2-2d'ac(\Delta\Delta-dd')
\end{aligned}
\]
und zwar die erste aus $12.15-13.13$, die zweite aus $14.14-12.17$; daraus sieht man leicht, dass notwendig $\Delta\Delta-dd'=0$ ist, ob nun $a=0$ oder nicht $=0$ gilt\HistoricalNote{*)}{Diese Herleitung der Gleichung $\Delta\Delta=dd'$ genügt für unseren gegenwärtigen Zweck; sonst könnten wir eine elegantere, hier aber zu lange Analyse geben, indem wir aus den Gleichungen $1\ldots11$ unmittelbar die Gleichung $0=(\Delta\Delta-dd')^2$ herleiten.}. Wir nehmen also an, dass in den Gleichungen $14,15,20,21$ auf der rechten Seite $\Delta\Delta-dd'$ gestrichen wird.

Setzen wir nun
\[
\begin{aligned}
\mathfrak{A}P+\mathfrak{B}(R-S)+\mathfrak{C}U&=mn'\\
\mathfrak{A}'Q+\mathfrak{B}'(R+S)+\mathfrak{C}'T&=m'n
\end{aligned}
\]
(wobei genau zu beachten ist, dass $n,n'$ auch Brüche sein können, obwohl $mn',m'n$ notwendig ganze Zahlen sind), so folgt leicht aus den Gleichungen $12\ldots17$
\[
Dmmn'n'=d'(\mathfrak{A}a+2\mathfrak{B}b+\mathfrak{C}c)^2=d'mm
\]
und ebenso aus den Gleichungen $18\ldots23$
\[
Dm'm'nn=d(\mathfrak{A}'a'+2\mathfrak{B}'b'+\mathfrak{C}'c')^2=dm'm'
\]

Es ist also $d=Dnn,\ d'=Dn'n'$, und wir erhalten den \textsc{ersten Schluss}: \textit{Die Determinanten der Formen $F,f,f'$ stehen notwendig im Verhältnis von Quadraten zueinander;} und den \textsc{zweiten}: \textit{$D$ teilt stets die Zahlen $dm'm',d'mm$.} Offenbar haben also $D,d,d'$ dasselbe Vorzeichen, und es kann keine Form in das Produkt $ff'$ transformierbar sein, deren Determinante größer ist als der größte gemeinsame Teiler der Zahlen $dm'm',d'mm$.

Multiplizieren wir die Gleichungen $12,13,14$ jeweils mit $\mathfrak{A},\mathfrak{B},\mathfrak{C}$; ebenso die Gleichungen $13,15,16$ und $14,16,17$ mit denselben Zahlen; addieren wir je drei Produkte und teilen die Summe durch $Dmn'$, wobei wir für $d'$ den Wert $Dn'n'$ schreiben. Dann ergibt sich
\[
P=an',\quad R-S=2bn',\quad U=cn'
\]

Multipliziert man ebenso die Gleichungen $18,19,20$ sowie $19,21,22$ und $20,22,23$ jeweils mit $\mathfrak{A}',\mathfrak{B}',\mathfrak{C}'$, so erhält man
\[
Q=a'n,\quad R+S=2b'n,\quad T=c'n
\]

\SourceLeaf{250}
Daraus folgt der \textsc{dritte Schluss}: \textit{Die Zahlen $a,2b,c$ sind den Zahlen $P,R-S,U$ proportional; setzt man ihr Verhältnis zu diesen gleich 1 zu $n'$, so ist $n'$ eine Quadratwurzel aus $\frac{d'}{D}$; ebenso haben die Zahlen $a',2b',c'$ zu $Q,R+S,T$ dasselbe Verhältnis, und setzt man dieses gleich 1 zu $n$, so ist $n$ eine Quadratwurzel aus $\frac{d}{D}$.}

Die Größen $n,n'$ können aber positive oder negative Quadratwurzeln aus $\frac{d}{D},\frac{d'}{D}$ sein. Daraus gewinnen wir eine Unterscheidung, die auf den ersten Blick unfruchtbar erscheinen mag, deren Nutzen sich aber im Folgenden hinreichend zeigen wird. Wir sagen nämlich, dass bei der Transformation der Form $F$ in $ff'$ die Form $f$ \textit{direkt} genommen wird, wenn $n$ positiv ist, und \textit{invers}, wenn $n$ negativ ist; ebenso wird $f'$ direkt oder invers genommen, je nachdem, ob $n'$ positiv oder negativ ist. Kommt die Bedingung hinzu, dass $k$ $=1$ ist, so heißt die Form $F$ \textit{entweder} aus beiden Formen $f,f'$ direkt zusammengesetzt, \textit{oder} aus beiden invers, \textit{oder} aus $f$ direkt und aus $f'$ invers, \textit{oder} aus $f$ invers und aus $f'$ direkt, je nachdem, ob $n,n'$ beide positiv, beide negativ, die erste positiv und die zweite negativ oder die erste negativ und die zweite positiv sind. Übrigens wird jeder leicht einsehen, dass diese Beziehungen nicht von der Reihenfolge der Formen $f,f'$ abhängen (siehe die erste Anmerkung zu diesem Artikel).

Weiter bemerken wir, dass der größte gemeinsame Teiler der Zahlen $P,\allowbreak Q,\allowbreak R,\allowbreak S,\allowbreak T,\allowbreak U$, nämlich $k$, die Zahlen $mn',m'n$ teilt (wie aus den oben festgestellten Werten hervorgeht), also das Quadrat $kk$ die Zahlen $mmn'n',m'm'nn$ und $Dkk$ die Zahlen $d'mm,dm'm'$ teilt. Umgekehrt teilt aber auch jeder gemeinsame Teiler von $mn',m'n$ die Zahl $k$. Sei nämlich $e$ ein solcher Teiler; er teilt offenbar auch die Zahlen $an',2bn',cn',a'n,2b'n,c'n$, \textit{d.h.}, die Zahlen $P,R-S,U,Q,R+S,T$ und damit auch $2R$ und $2S$. Wäre nun $\frac{2R}{e}$ ungerade, so müsste auch $\frac{2S}{e}$ ungerade sein (weil Summe und Differenz gerade sind), also auch ihr Produkt. Dieses Produkt ist aber $=\frac{4}{ee}(b'b'nn-bbn'n')=\frac{4}{ee}(d'nn+a'c'nn-dn'n'-acn'n')=\frac{4}{ee}(a'c'nn-acn'n')$ und damit gerade, weil $e$ die Zahlen $a'n,c'n,an',cn'$ teilt. Also ist $\frac{2R}{e}$ notwendig gerade und damit $R$ sowie $S$ durch $e$ teilbar. Da also $e$ alle sechs Zahlen $P,Q,R,S,T,U$ teilt, teilt es auch deren größten gemeinsamen Teiler $k$. \textit{Q. E. D.} --- Daraus folgt, dass $k$ der größte gemeinsame Teiler der Zahlen $mn',m'n$ ist; daraus sieht man leicht, dass \textit{$Dkk$ der größte gemeinsame Teiler der Zahlen $dm'm',d'mm$ ist.} Das ist der \textsc{vierte Schluss}. Offenbar ist also, wenn $F$ aus $f$ und $f'$ zusammengesetzt ist, $D$ der größte gemeinsame Teiler der Zahlen $dm'n',d'mm$, und umgekehrt; diese Eigenschaft hätte man auch als Definition der zusammengesetzten Form nehmen können. Eine aus den
\SourceLeaf{251}
Formen $f,f'$ zusammengesetzte Form hat also die größtmögliche Determinante unter allen Formen, die in das Produkt $ff'$ transformierbar sind.

Bevor wir weitergehen können, müssen wir zunächst den Wert von $\Delta$ genauer bestimmen. Wir haben zwar gezeigt, dass er $=\surd dd'=\surd DDnnn'n'$ ist, doch sein \textit{Vorzeichen} ist dadurch noch nicht bestimmt. Dazu leiten wir aus den Fundamentalgleichungen 1---11 die Gleichung $DPQ=\Delta aa'$ her (sie ergibt sich aus $5.6-1.11$), also $Daa'nn'=\Delta aa'$, woraus, wenn nicht eine der Zahlen $a,a'$ $=0$ ist, $\Delta=Dnn'$ folgt. Ganz ebenso lassen sich aber aus den Fundamentalgleichungen acht weitere herleiten, in denen auf der linken Seite $Dnn'$ und auf der rechten $\Delta$ mit $2ab',ac',2ba',4bb',2bc',ca',2cb',cc'$ multipliziert sind\HistoricalNote{*)}{Die Analyse, die die Leser leicht selbst finden können, müssen wir der Kürze wegen weglassen.}. Da weder alle $a,2b,c$ noch alle $a',2b',c'$ $=0$ sein können, folgt leicht, dass in allen Fällen $\Delta=Dnn'$ gilt und damit $\Delta$ dasselbe Vorzeichen wie $D,d,d'$ oder das entgegengesetzte hat, je nachdem, ob $n,n'$ gleiche oder verschiedene Vorzeichen haben.

Weiter bemerken wir, dass die Zahlen $aa',\allowbreak 2ab',\allowbreak ac',\allowbreak 2ba',\allowbreak 4bb',\allowbreak 2bc',\allowbreak ca',\allowbreak 2cb',\allowbreak cc',\allowbreak 2bb'+2\Delta,\allowbreak 2bb'-2\Delta$ alle durch $mm'$ teilbar sind. Für die ersten neun ist das unmittelbar klar; für die beiden übrigen lässt es sich ebenso beweisen, wie wir oben gezeigt haben, dass $R$ und $S$ durch $e$ teilbar sind. Offenbar sind nämlich $4bb'+4\Delta$ und $4bb'-4\Delta$ durch $mm'$ teilbar (weil $4\Delta=\surd16dd'$ und $4d$ durch $mm$, $4d'$ durch $m'm'$ teilbar ist, also $16dd'$ durch $mmm'm'$ und $4\Delta$ durch $mm'$), und die Differenz der Quotienten ist gerade; ihr Produkt lässt sich leicht als gerade nachweisen, also ist jeder Quotient gerade und $2bb'+2\Delta,2bb'-2\Delta$ sind durch $mm'$ teilbar.

Aus den elf Fundamentalgleichungen folgen nun leicht die sechs folgenden:
\[
\begin{aligned}
APP&=aa'q'q'-2ab'qq'+ac'qq\\
AQQ&=aa'q''q''-2ba'qq''+ca'qq\\
ARR&=aa'q'''q'''-2(bb'+\Delta)qq'''+cc'qq\\
ASS&=ac'q''q''-2(bb'-\Delta)q'q''+ca'q'q'\\
ATT&=ac'q'''q'''-2bc'q'q'''+cc'q'q'\\
AUU&=ca'q'''q'''-2cb'q''q'''+cc'q''q''
\end{aligned}
\]

Daraus folgt, dass alle $APP,AQQ$ usw. durch $mm'$ teilbar sind; da $kk$ der größte gemeinsame Teiler der Zahlen $PP,QQ,$
\SourceLeaf{252}
$RR$ usw. ist, folgt daraus leicht, dass auch $Akk$ durch $mm'$ teilbar ist. Setzt man für $a,2b,c$, $a',2b',c'$ ihre Werte $\frac{P}{n'}$ usw. oder $\frac{1}{n'}(pq'-qp')$ usw. ein, so gehen sie in sechs andere Gleichungen über, auf deren rechten Seiten die Produkte aus der Größe $\frac{1}{nn'}(q'q''-qq''')$ mit $PP,QQ,RR$ usw. stehen. Diese sehr leichte Rechnung überlassen wir den Lesern. Daraus folgt (weil nicht alle $PP,QQ$ usw. $=0$ sein können) $Ann'=q'q''-qq'''$.

Ebenso leitet man aus den Fundamentalgleichungen sechs weitere Gleichungen her, die sich von den vorigen nur dadurch unterscheiden, dass überall $A$ durch $C$ und $q,q',q'',q'''$ jeweils durch $p,p',p'',p'''$ ersetzt werden. Der Kürze wegen schreiben wir sie nicht hin. Daraus folgt ebenso, dass $Ckk$ durch $mm'$ teilbar ist und $Cnn'=p'p''-pp'''$ gilt.

Schließlich ergeben sich aus derselben Quelle diese sechs Gleichungen:
\[
\begin{aligned}
BPP&=-aa'p'q'+ab'(pq'+qp')-ac'pq\\
BQQ&=-aa'p''q''+ba'(pq''+qp'')-ca'pq\\
BRR&=-aa'p'''q'''+(bb'+\Delta)(pq'''+qp''')-cc'pq\\
BSS&=-ac'p''q''+(bb'-\Delta)(p'q''+q'p'')-ca'p'q'\\
BTT&=-ac'p'''q'''+bc'(p'q'''+q'p''')-cc'p'q'\\
BUU&=-ca'p'''q'''+cb'(p''q'''+q''p''')-cc'p''q''
\end{aligned}
\]
woraus wir wie zuvor schließen, dass $2Bkk$ durch $mm'$ teilbar ist und $2Bnn'=pq'''+qp'''-p'q''-q'p''$ gilt.

Da also $Akk,2Bkk,Ckk$ durch $mm'$ teilbar sind, sieht man leicht, dass auch $Mkk$ durch $mm'$ teilbar sein muss. Aus den Fundamentalgleichungen folgt aber, dass $M$ die Zahlen $aa',\allowbreak 2ab',\allowbreak ac',\allowbreak 2ba',\allowbreak 4bb',\allowbreak 2bc',\allowbreak ca',\allowbreak 2cb',\allowbreak cc'$ teilt und damit auch $am',2bm',cm'$ (die größten gemeinsamen Teiler jeweils der ersten, mittleren und letzten drei); schließlich auch $mm'$, ihren größten gemeinsamen Teiler. Daraus sieht man, dass in dem Fall, in dem die Form $F$ aus den Formen $f,f'$ zusammengesetzt ist, also $k=1$ gilt, notwendig $M=mm'$ ist. Das ist der \textsc{fünfte Schluss}.

Ist der größte gemeinsame Teiler der Zahlen $A,B,C$ gleich $\mathfrak{M}$, so ist dieser entweder $=M$ (wenn die Form $F$ eigentlich primitiv oder aus einer eigentlich primitiven Form abgeleitet ist) oder $=\frac{1}{2}M$ (wenn $F$ uneigentlich primitiv oder aus einer uneigentlich primitiven Form abgeleitet ist); ebenso bezeichnen wir die größten gemeinsamen Teiler der Zahlen $a,b,c;\ a',b',c'$ jeweils mit $\mathfrak{m},\mathfrak{m}'$; dann ist $\mathfrak{m}$ entweder $=m$ oder $\frac{1}{2}m$ und $\mathfrak{m}'$ entweder $=m'$ oder $=\frac{1}{2}m'$. Offenbar teilt nun $\mathfrak{m}\mathfrak{m}$ die Zahl $d'$, $\mathfrak{m}'\mathfrak{m}'$ die Zahl $d'$, also $\mathfrak{m}\mathfrak{m}\mathfrak{m}'\mathfrak{m}'$ die Zahl $dd'$ oder $\Delta\Delta$ und $\mathfrak{m}\mathfrak{m}'$ die Zahl $\Delta$.
\SourceLeaf{253}
Aus den letzten sechs Gleichungen für $BPP$ usw. folgt dann, dass $\mathfrak{m}\mathfrak{m}'$ die Zahl $Bkk$ teilt und damit (weil es auch $Akk,Ckk$ teilt) auch $\mathfrak{M}kk$. Wenn also $F$ aus $f,f'$ zusammengesetzt ist, teilt $\mathfrak{m}\mathfrak{m}'$ die Zahl $\mathfrak{M}$. Sind daher in diesem Fall beide $f,f'$ eigentlich primitiv oder aus eigentlich primitiven Formen abgeleitet, gilt also $\mathfrak{m}\mathfrak{m}'=mm'=M$, so ist $\mathfrak{M}=M$, oder $F$ ist eine ebensolche Form. Sind dagegen unter derselben Voraussetzung entweder beide $f,f'$ oder wenigstens eine davon uneigentlich primitiv oder aus einer uneigentlich primitiven Form abgeleitet, \textit{zum Beispiel} die Form $f$, so folgt aus den Fundamentalgleichungen, dass $aa',2ab',ac',ba',2bb',bc',ca',2cb',cc'$ durch $\mathfrak{M}$ teilbar sind, also auch $am',bm',cm'$ und damit auch $\mathfrak{m}m'=\frac{1}{2}mm'=\frac{1}{2}M$; daraus folgt in diesem Fall notwendig $\mathfrak{M}=\frac{1}{2}M$, oder die Form $F$ ist ebenfalls uneigentlich primitiv oder aus einer uneigentlich primitiven Form abgeleitet. Das ergibt den \textsc{sechsten Schluss}.

Schließlich bemerken wir: Wenn man \textit{annimmt}, dass die neun Gleichungen
\[
\begin{gathered}
an'=P,\quad 2bn'=R-S,\quad cn'=U\\
a'n=Q,\quad 2b'n=R+S,\quad c'n=T\\
\begin{gathered}Ann'=q'q''-qq''',\\2Bnn'=pq'''+qp'''-p'q''-q'p'',\\Cnn'=p'p''-pp'''\end{gathered}
\end{gathered}
\]
(die wir mit $\Omega$ bezeichnen werden, weil wir im Folgenden öfter darauf zurückkommen müssen) gelten, und dabei sogar $n,n'$ als Unbekannte betrachtet, von denen jedoch keine $=0$ ist, dann bestätigt man durch Einsetzen leicht, dass auch die Fundamentalgleichungen $1\ldots9$ notwendig wahr sind, also die Form $(A,B,C)$ durch die Substitution $p,p',p'',p''';\ q,q',q'',q'''$ in das Produkt der Formen $(a,b,c)(a',b',c')$ übergeht; außerdem gilt
\[
bb-ac=nn(BB-AC),\quad b'b'-a'c'=n'n'(BB-AC)
\]
Die Rechnung, die hier zu lang wäre, überlassen wir dem Fleiß der Leser.
